\section{Zadanie 12}
Niech \(X_1, X_2, \ldots, X_n\) będą niezależnymi zmiennymi losowymi z tego samego rozkładu o skończonej wartości oczekiwanej \(\mu\)
oraz skonczonym odchyleniu standardowym \(\sigma\). Na bazie nierówności Czebyszewa pokaż, że dla kazdego \(\epsilon > 0\)

\[\lim_{n \to \infty} P\left(\left|\frac{X_1 + X_2 + \ldots + X_n}{n} - \mu\right| > \epsilon\right) = 0\]

\subsection*{Rozwiązanie}
Zauważmy, że \(\frac{X_1 + X_2 + \ldots + X_n}{n}\) jest średnią arytmetyczną zmiennych losowych \(X_i\). Zastosujmy nierówność Czebyszewa:

\[P\left(\left|\frac{X_1 + X_2 + \ldots + X_n}{n} - \mu\right| > \epsilon\right) \leq \frac{\variance{\frac{X_1 + X_2 + \ldots + X_n}{n}}}{\epsilon^2}\]

Obliczmy wariancję średniej:

\[\variance{\frac{X_1 + X_2 + \ldots + X_n}{n}} = \frac{1}{n^2} \sum_{i=1}^n \variance{X_i} = \frac{\sigma^2}{n}\]

Podstawiając to do nierówności Czebyszewa, otrzymujemy:

\[P\left(\left|\frac{X_1 + X_2 + \ldots + X_n}{n} - \mu\right| > \epsilon\right) \leq \frac{\sigma^2/n}{\epsilon^2} = \frac{\sigma^2}{n\epsilon^2}\]

Zauważmy, że gdy \(n \to \infty\), to \(\frac{\sigma^2}{n\epsilon^2} \to 0\). Zatem:

\[\lim_{n \to \infty} P\left(\left|\frac{X_1 + X_2 + \ldots + X_n}{n} - \mu\right| > \epsilon\right) = 0\]

\subsection*{Odpowiedź}
\[
	\boxed{\text{Tak to prawda}}
\]
